% AI-written, AI cross-reviewed; no human sign-off; not a claim that RH is proved
% Atlas unit D351. Local manuscript v3, item G.364, heading on printed PDF p. 247.
% ReciprocityTheory.AnForm.M.poincareLelong
% paper2/decomposition_tables.tex:3014–3019. Verbatim excerpt; SHA-256 5da7d86300d4ceb44e36e8a9a27e0fa6828a0cb82bb06c1510d5f0ffd587b2bb.
% Original TeX uses the macros and cross-reference keys of the source manuscript.


% Start of atlas edition annotation. The following block is outside the original excerpt.
\paragraph{Status of named inputs in this reading.}
\noindent\nolinkurl{ReciprocityTheory.AnForm.M.poincareLelong} --- in-repo constant (\path{ReciprocityTheory/AnalyticForms/PotentialTheory/PLLocalStatement.lean:27}); not compared.\par
% End of atlas edition annotation.

\begin{formalstatement}[D351]\label{decomp:D351}
On a compact Hausdorff space with a finite holomorphic atlas of complex dimension e+1, let a section of a metrized line datum have a simple normal crossings divisor with multiplicities \nolinkurl{m_k} and holomorphic component inclusions. For any smooth test function and e admissible potentials, integration against the log norm equals the sum of the component integrals weighted by \nolinkurl{m_k} minus the ambient curvature integral.
\end{formalstatement}
{\small\raggedright\noindent\textsf{\bfseries Related reference}\quad \nolinkurl{ReciprocityTheory.AnForm.M.poincareLelong}\par\noindent \nolinkurl{ReciprocityTheory/AnalyticForms/PotentialTheory/PLLocalStatement.lean:27--27}\par\noindent SHA-256 \nolinkurl{c088e29fe61e4ddb419f38ea664da4f5ab313e1cf1c590ec2e4036a43eafe8e0}\par\noindent Dominated 975\quad Disjoint 975\quad terminal.\par}
The logarithmic current and the given normalization of the curvature and divisor must be identified.
\hyperref[lit:S031]{Published sources and search S031}.
